% 本文件由 LaTeX 排版 Agent 根据有效 translation.json 自动生成。
% author=Tertullian; work_id=0318; title=蝎子毒

\chapter{蝎子毒}

% structure: p0001 source=h1 target=omitted reason=page-title-already-rendered-by-parent

\Emph{蝎子蜇伤的解毒剂。}\par

% structure: p0003 source=h2 target=section reason=relative-heading-rank; base=h2

\section{第一章}

大地仿佛从脓疮中生出渺小蝎子所带来的巨大祸害。其毒物与种类一样多，其灾祸与类别一样多，其疼痛与颜色一样多。尼坎德记载了\Emph{关于蝎子}，并描绘了它们。然而，它们全都用尾巴（这尾巴是任何从身体后端伸出的东西，并作鞭打）发动攻击时使用同一种动作。因此，蝎子体内那串节结，在内部是一条细微的毒脉，像弓弦一样弹起，末端拉紧一根带倒刺的毒针，如同投射弹丸的机械。由此，人们也将一种战争器械称为蝎子，它在被拉回后给予箭矢以动力。它们的针尖也是极细的导管，用以造成伤口；刺入时便注入毒液。通常的危险季节是夏季：当南风和西南风起时，凶暴便扬帆。在治疗方面，自然所提供的\Emph{某些}物质非常有效；巫术也缠上绷带；医术用柳叶刀和吸杯来对抗。有些人急忙预先饮用保护剂；但性交会消耗它，他们又变得干燥。我们有信德作为防御，如果我们不立即画十字、驱魔并以那野兽涂抹脚跟，那么连不信本身也不会击中我们。最后，我们也经常以此帮助外邦人，因为我们已从天主获得了那宗徒先运用过的能力，他曾藐视毒蛇的咬伤。\BibleRef{宗 28:3}那么，这支笔提供什么呢？如果信德靠自身所拥有的就已安全？为了使它也能在其它时候靠自身所拥有的安全，当它遭受自身的蝎子时。这些蝎子也有麻烦的小体型，种类不同，都用一种方式武装，在特定时间被激怒，而且不是别的时刻，正是酷热之时。这在基督徒中是\Emph{迫害}的季节。因此，当信德大受震荡，教会燃烧，正如荆棘丛所代表的，\BibleRef{出 3:2}那时诺斯替派爆发，瓦伦提尼派爬行出来，所有殉道的反对者涌出，他们自己也热切地想刺、穿透、杀死。因为他们知道许多人天真、无经验，且软弱，事实上极大多数基督徒随风摇摆、顺从情绪，他们意识到只有在恐惧打开了灵魂的入口时，尤其是在某种\Emph{残暴的展示}已经以殉道者的信德戴上冠冕时，他们才最可接近。因此，他们一直拖着尾巴，首先将它触及情感，或者像在空地上挥舞鞭子。无辜者遭受如此苦难。所以你会以为\Emph{说话者}是一个弟兄或一位较好的外邦人。一个不为难任何人的教派就这样被对待！然后他们刺入。人毫无道理地丧命。因为\Quote{他们丧命，毫无道理}是第一次插入。然后他们现在致命地击打。但天真的灵魂不知道所写的是什么，以及它有什么含义，我们应当在何处、何时、在谁面前承认（\Emph{或应当}），只知道为天主而死，既然祂保存我，那甚至不是天真，而是愚昧，不，是疯狂。如果祂杀死我，祂怎么有义务保存我？基督一次为我们死了，一次被击杀，以便我们不被击杀。如果祂向我要求同样的回报，祂也期望从我暴力死亡中获得救恩吗？或者天主渴望人的血，尤其是当祂拒绝公山羊和公山羊的血时？的确，祂更愿罪人悔改而不是死亡。\BibleRef{则 33:11}而祂如何渴望那些不是罪人的人的死亡？这些以及或许其他含有异端毒物的狡诈伎俩，谁会不被它们刺入，或是为了疑惑（若不是为了毁灭），或是为了激怒（若不是为了死亡）？因此，至于你，如果信德警觉，请立即用你的鞋底尽可能诅咒蝎子，让它昏迷死去？但如果它饱饮伤口，便将毒液向内驱赶，使之加速进入内脏；立刻所有先前的知觉变得迟钝，心灵的血液冻结，灵魂的肉体憔悴，对基督徒之名的厌恶伴随着酸味。理解力也为自己寻找呕吐之处；于是，一度被击伤的软弱在异端或异教中呼出受伤的信德。现在的情况是\Emph{如此，我们正身处}极度酷热中，迫害的犬星——无疑起源于那狗头者本人。一些基督徒被火、一些被剑、一些被野兽试验；其他人因被棍棒和爪钩折磨而暂时尝到殉道滋味，在监狱中挨饿。我们自己被指定追逐，如同兔子从远处被包围；而异端者照常四处走动。因此，时局促使我用笔准备针对扰乱我们教派的小兽的解毒剂，以便我得以治愈。你读的时候同时喝。这药汁并不苦。如果主的话语比蜂蜜和蜂巢更甜，汁液便源于此。如果天主的应许流溢奶与蜜，\BibleRef{出 3:17}构成那药汁的成分便带有这味道。\Quote{但祸哉，那些把甜变为苦，把光明变为黑暗的人。}\BibleRef{依 5:20}因为同样，那些反对殉道、把救恩说成毁灭的人，也将甜变为苦，把光明变为黑暗；这样，他们偏爱这极其悲惨的生命胜过那最蒙福的生命，就把苦当作甜，把黑暗当作光明。\par

% structure: p0005 source=h2 target=section reason=relative-heading-rank; base=h2

\section{第二章}

但是，关于从殉道中获得的好处，我们还不能学习，除非我们首先\Emph{听}到受苦的本分；我们也不能\Emph{学习}它的益处，除非我们先\Emph{听}到它的必要性。天主的许可（问题）在先——即天主是否愿意并命令了这类事情，以便那些断言它不好的异端者，只有在被制服之后，才被提供理由认为它有益。合适的是，异端者应被驱向本分，而不是被引诱。顽固必须被征服，而不是被哄骗。而且，当然，那将被判定为足够好的，正是那被证明是由天主所设立并命令的。让福音稍等片刻，让我先陈述它们的根源——法律，让我从那些我藉以想起天主本身的经文中，确认天主的旨意：祂说：\Quote{\Emph{我}是，}祂说，\Quote{天主，你的天主，曾将你从埃及地领出来。除我以外，你不可有别的神。你不可为自己制造偶像，也不可制造天上、地下、水中之物的任何形象。你不可崇拜它们，也不可事奉它们。因为我是上主你的天主。}\BibleRef{出 20:2}\newline{}同样，在同一部\BibleBook{出谷}{纪}中：\Quote{你们亲自看见我从天上与你们说话。你们不可制造银神，也不可制造金神。}\BibleRef{出 20:22}\newline{}在\BibleBook{申命}{纪}中也有如下的话：\Quote{以色列，请听：上主你的天主是唯一的上主；你当全心、全意、全力爱上主你的天主。}\BibleRef{申 6:4}\newline{}又说：\Quote{不要忘记上主你的天主，祂曾将你从埃及地、从为奴之家领出来。你要敬畏上主你的天主，惟独事奉祂，依附祂，以祂的名起誓。不可追随别的神，即你们四周列国的神，因为上主你的天主在你们中间是忌邪的天主，免得祂的怒气向你们发作，将你们从地上消灭。}\BibleRef{申 6:12}\newline{}但在向他们提出祝福和诅咒时，祂也说：\Quote{如果你们听从上主你的天主的诫命，就是我今日所吩咐你们的，不偏离我所吩咐你们的道路，去事奉你们不认识的神，祝福就必临到你们。}\BibleRef{申 11:27}\newline{}至于彻底铲除它们：\Quote{你们要彻底毁坏你们要继承的列国事奉他们神明的所有地方，无论在高山、丘陵或绿树下。你们要拆毁他们的祭坛，打碎他们的柱像，砍下他们的木偶，用火焚烧他们神明的雕像，并从那地方除去他们的名号。}\BibleRef{申 12:2}\newline{}祂进一步敦促，当他们（以色列人）进入应许之地并赶出那里的列国时：\Quote{你要谨慎，不可在他们从你面前被赶出之后追随他们，也不可求问他们的神，说：\InnerQuote{列国怎样事奉他们的神，我也要照样行。}}\BibleRef{申 12:30}\newline{}但祂也说：\Quote{如果你们中间有先知或作梦的兴起，给你一个征兆或奇事，这征兆或奇事应验了，他却说：\InnerQuote{我们去事奉别的神吧}，就是你们不认识的神，你们不要听信那先知或作梦的话，因为上主你们的天主正在考验你们，要知道你们是否全心全力爱祂。你们要随从上主你们的天主，敬畏祂，遵守祂的诫命，听从祂的声音，事奉祂，依附祂。但那先知或作梦的必须处死，因为他唆使你们背离上主你们的天主。}\BibleRef{申 13:1}\newline{}但在另一处又说：\Quote{然而，如果你的弟兄，你父亲或母亲的儿子，或你的儿子，或你的女儿，或你怀中的妻子，或如同你灵魂的朋友，暗中诱骗你说：\InnerQuote{我们去事奉别的神吧}，就是你和你的祖先都不认识的神，是你们周围或远或近的列国的神，你不可随从他，也不可听从他。你的眼不可顾惜他，不可怜悯他，也不可掩护他；你必须告发他。你的手要先加在他身上将他处死，然后众民的手也加在他身上。你要用石头砸死他，因为他企图使你背离上主你的天主。}\BibleRef{申 13:6}\newline{}祂同样关于城市补充说，如果发现有城市因不义之人的怂恿而转向别的神，那么城中所有居民都要被杀死，一切所有都要成为禁物，所有战利品都要聚集到城中各出口，连同所有百姓，在街头用火烧尽，在上主天主的眼前；并且祂说：\Quote{它永远不可再被居住；它不可再被建造；凡从这禁物中取走的，都不可粘在你的手上，好使上主转消祂的烈怒。}\BibleRef{申 13:16}\newline{}因祂憎恶偶像，也制定了一系列诅咒：\Quote{凡制造雕像或铸像，即匠人手所做的可憎之物，并把它放在隐秘处的，是可诅咒的。}\BibleRef{申 27:15}\newline{}但在\BibleBook{肋未}{纪}中祂说：\Quote{不可转向偶像，也不可为自己铸造神像：我是上主你们的天主。}\BibleRef{默 19:4}\newline{}在别处又说：\Quote{以色列子民是我的家仆；是我将他们从埃及地领出来的：我是上主你们的天主。你们不可为自己制造手造的偶像，也不可竖立雕像，也不可在你们的地上树立一块引人注目的石头（来崇拜它）：我是上主你们的天主。}这些话最初是由上主藉梅瑟的口所说的，当然适用于任何像这样被上主以色列的天主从一个最迷信的世界——埃及，并从人类奴役之地领出来的人。但从每一位相继的先知口中，也传出同一位天主的话语，以重新颁布同样的命令来加强祂同样的法律，并且首先特别宣布，没有其他本分比警惕一切制造和崇拜偶像更为重要；正如祂藉达味的口所说：\Quote{列国的神是金银：它们有眼，却看不见；有耳，却听不见；有鼻，却闻不到；有口，却不能说话；有手，却不能触摸；有脚，却不能行走。制造它们并依靠它们的，必像它们一样。}\par

% structure: p0007 source=h2 target=section reason=relative-heading-rank; base=h2

\section{第三章}

我也不认为有必要讨论天主采取一个可敬的途径，禁止祂自己的名号与尊荣被交于谎言，或是不允许那些祂从假宗教的迷宫中夺回的人再次返回埃及，或是不让祂所拣选的人离开祂。因此，我们也不必讨论：祂是否愿意遵守祂所选定设立的法则，以及祂是否公正地惩罚对祂希望遵守之法则的废弃；因为如果祂不希望它被遵守，祂的设立就是徒劳的；如果祂不愿意维护它，祂希望它被遵守也是徒劳的。的确，我的下一步是检验天主反对假宗教的这些设立——那些完全被击败的以及受惩罚的，因为整个关于殉道的论证都将依赖于此。梅瑟独自与天主在山上，百姓不能容忍他如此必要的缺席，便为自己寻求制造神祇，而他本人则将更喜欢毁灭它们。\EditorialNote{出谷纪第三十二章} 亚郎被催逼，吩咐将妇女的耳环收集起来，投入火中。因为百姓即将作为对自己的惩罚，失去真正的耳饰——天主的话语。智慧的火为他们铸成了一个牛犊的像，责备他们心之所向即财宝所在——即在埃及，那里如同其他动物一样，一头牛也被视为神圣。因此，他们的近亲杀了三千人，因为他们冒犯了如此亲近的天主，这庄严地标示了过犯的开端与应得之罚。正如我们在\WorkTitle{户籍纪}中读到，\BibleRef{户 25:1}以色列在史廷偏离正道，百姓去与摩阿布的女儿行淫，满足情欲：他们被偶像引诱，以致也在心灵上犯奸淫；最后，他们吃了玷污的\Emph{祭献}；然后他们既崇拜那民族的神祇，又被接纳参加贝耳培敖尔的礼仪。对于这次堕入与奸淫为姐妹的偶像崇拜，他们被同胞的刀剑杀死二万三千人，才平息了天主的义怒。若苏厄·农之子死后，他们离弃了他们祖先的天主，事奉偶像——巴耳诸神和阿市托勒特；\BibleRef{民 2:8} 主发怒，将他们交在劫掠者手中，他们继续被他们劫掠，被卖予仇敌，在敌人面前完全不能站立。无论他们往哪里去，天主的手降下灾祸攻击他们，他们极其困苦。此后，天主在他们中间设立\OriginalTerm{判官}{critas}，如同我们的监察官。但他们甚至对这些判官也不恒心服从。每当一位判官去世，他们就比他们的祖先更甚地行恶，去追随别的神祇，事奉并敬拜它们。因此主发怒。\Quote{的确，}祂说，\Quote{这民族违背了我与他们祖先订立的盟约，不听我的声音，我也必不理会，不再从他们面前除去约书亚死时所留下的各国中的人。}\BibleRef{民 2:20} 因此，在几乎所有判官及继任国王的年鉴中，四周列国的力量得以保存，每当他们离弃祂，尤其是转向偶像崇拜时，祂就以战争、被掳和外族的轭量给以色列愤怒。\par

% structure: p0009 source=h2 target=section reason=relative-heading-rank; base=h2

\section{第四章}

因此，既然从一开始这种敬拜就已被禁止——以那些众多而重要的诫命为证——并且从未在未随后受罚的情况下进行过，正如如此众多而令人印象深刻的例子所示，并且没有任何冒犯被天主视为如此僭越，如同这种侵犯一样，我们应当进一步领悟神圣威胁及其应验的意旨，这甚至在当时不仅通过不加质疑，还通过忍受殉道而受到了称赞，祂无疑是借禁止偶像崇拜而提供了殉道的机会。否则，殉道就不会发生。而且，祂当然为此提供了祂自己的权威作为保证，愿意那些祂曾给予机会让其发生的事件得以成就。如今（这很重要），因为我们正受到关于天主旨意的严厉刺扎，而蝎子不断重复其螫刺，否认这个旨意的存在，挑剔它，以致他要么暗示有另一位神，以至于这不是他的旨意，要么无论如何推翻我们的神，既然他的旨意如此，要么彻底否认天主\Emph{这个}旨意，如果他不能否认天主本身的话。但是，就我们而言，我们在别处论述天主以及异端教导的全部其余部分，现在我们为一种交锋划定明确的界线，坚持认为这个旨意——就是导致殉道的那个——并非属于另一位神，而正是以色列的天主，其根据是关于始终被禁止的偶像崇拜的诫命，以及对受惩罚的偶像崇拜的审判。因为如果遵守诫命涉及遭受暴力，那么这可以说就是一个关于遵守诫命的诫命，要求我承受那使我能够遵守诫命的事物，也就是暴力，无论何种暴力，当我在提防偶像崇拜时威胁我。并且当然（在所假设的情形中），诫命的颁布者强制人服从它。因此，祂不可能不愿意那些事件发生，通过它们，这服从将得以显明。我领受的禁令是：不可提及任何别的神，甚至不可用言语——用舌头或用手——塑造一个神，并且不可敬拜或以任何方式向别的神表示敬意，唯独向那如此命令我的祂，我既被命令敬畏祂，以不被祂遗弃，又要用我整个存在爱祂，好为祂而死。在这誓言之下服兵役，我受到敌人的挑战。如果我向他们投降，我就和他们一样。持守这誓言，我在战斗中奋勇作战，受伤，被砍成碎片，被杀死。谁希望他的士兵落得如此致命的结局呢？岂非那用这样的誓言印证他的人？\par

% structure: p0011 source=h2 target=section reason=relative-heading-rank; base=h2

\section{第五章}

因此，你有了我天主的旨意。我们已治愈了这一刺。让我们好好注意另一个关于祂旨意特质的打击。要表明我的天主是善的——这是马西翁派如今已由我们得知的真理——是冗长的。同时，既然祂被称为天主，就必须相信祂是善的，这便足够了。因为若有人假设天主是恶的，他便无法立足于这两个构成要素：他必得肯定他以为恶的那位不是天主，或他宣称是天主的那位是善的。因此，祂的旨意也将是善的，因为除非是善的，祂便不是天主。天主所愿意的事物本身——我指的是殉道——的善性也将表明这一点，因为唯有善者才愿意善事。我坚决主张殉道是善的，因为它是由那位同样禁止并惩罚偶像崇拜的天主所要求的。因为殉道奋力对抗并反对偶像崇拜。但奋力对抗并反对恶事不能不是善的。并非我否认恶事之间也有竞争，正如善事亦然；但这理由需要不同的状况。因为殉道与偶像崇拜竞争，并非出于它们共有的某种恶意，而是出于它自身的仁慈；因为它从偶像崇拜中解救出来。谁不会宣称那从偶像崇拜中解救出来的事是善的呢？偶像崇拜与殉道之间的对立，除了生与死，还能是什么呢？生命将被视为殉道，如同偶像崇拜被视为死亡。谁若称生命为恶，就必得说死亡为善。这种悖逆也属于人类——抛弃有益之物，接受有害之物，逃避一切危险的疗法，简言之，急于求死而非得医治。因为许多人逃避药物的帮助，许多人出于愚蠢，许多人出于恐惧和虚假的羞耻。医术显然有一种表面上的残酷，由于手术刀、烙铁和芥末的强烈热力；但被切割、烧灼、拉扯和咬啮，并不因此就是恶，因为它引起有益的疼痛；并且不会仅仅因为它带来痛苦就被拒绝，而是因为它不可避免地带来痛苦才被施用。那善的后果是为可怕的工作道歉。简言之，那个在医生手中嚎叫、呻吟、咆哮的人，不久就会将谢礼放在同一双手上，宣称他们是最好的操刀者，不再说他们残忍。同样，殉道也猛烈地发作，但为了救恩。天主也有自由藉由火与剑以及一切痛苦之事来医治以得永生。但至少在那方面你也会钦佩医生，因为他在治疗中多半使用与疾病性质相似的特性，当他以错误的方式帮助、通过那些导致痛苦的事物来救助时。他既以更大的热敷来以热制热；又通过不满足口渴、反而折磨来抑制炎症；并以每剂苦药收敛过多的胆汁，并再开一个小静脉来止血。但你会认为天主必须受责备，而且是因为嫉妒，如果祂选择与疾病抗争，并通过仿效疾病来行善，以死亡消灭死亡，以杀戮驱散杀戮，以折磨驱散折磨，以惩罚驱散惩罚，以收回生命赐予生命，以伤害肉体帮助肉体，以夺走灵魂保存灵魂。你所认为的悖逆是理性；你所认为的残忍是仁慈。因此，看到天主藉由短暂的痛苦成就永恒的医治，为自己的繁荣赞美你的天主吧；你已落入祂的手中，但幸运地落入了。祂也落入了你的疾病中。人总是首先为医生提供工作；简言之，他为自己带来了死亡的危险。他从自己的主那里，如同从医生那里，接受了足够有益的规则：按照法律生活，即可以吃园中所有的（果实），但只应戒绝一棵小树，当时医生自己知道那棵树是危险的。他听从了他偏爱的那一位，打破了自制。他吃了禁果，因过犯而饱足，遭受了致死的消化不良；他当然完全配得上失去全部生命，因为他想这样做。但因过犯引起的炎症肿瘤一直忍受，直到时机成熟时药剂被调和，主逐渐准备了医治的方法——信仰的一切规条，它们也带有与疾病原因相似的特征，因为它们以生命之言取消死亡之言，并以顺从的倾听减少过犯的倾听。因此，即使那位医生命令人死去，祂也驱散了死亡的昏睡。为什么人现在对治疗带来的痛苦感到不情愿，而那时对疾病带来的痛苦却不感到不情愿呢？难道他憎恨为救恩而死，而那时却不憎恨为毁灭而死？——对解毒剂感到恶心，而对毒药却张大了口？\par

% structure: p0013 source=h2 target=section reason=relative-heading-rank; base=h2

\section{第六章}

但若天主为竞赛之故，为我们指定了殉道，好使我们借此与对手较量，使祂继续捶打那人——人类曾甘愿受其捶打——那么在此，天主仍以宽仁而非严苛为怀。因为祂希望人——现已藉信德从魔鬼喉咙中被夺回——同样以勇气践踏他，使他不只逃脱，更彻底击败他的仇敌。那召叫他们得救恩的，也乐意召叫他们得光荣，好使那些因获得解救而欢欣的人，在同样被加冕时欣喜若狂。世界以何等善意庆祝那些竞赛，即希腊人的战斗节庆和迷信赛事，兼有敬拜与娱乐形式，如今在\Place{阿非利加洲}也变得明显。迄今，各城市分别发送贺词，困扰着\Place{迦太基}——它在竞技场年迈之后获得了\OriginalTerm{皮提亚竞技会}{Pythian game}。因此，世人认为以下是最恰当的方法以检验学业的熟练程度：将技艺形式置于竞争之中，引出身体与声音的现有状况，其奖赏为告密者，公开展示为裁判，愉悦为判决。凡有竞赛之处，必有创伤：拳头使人踉跄，脚跟像公羊般抵撞，拳击手套撕裂皮肉，鞭子留下深痕。然而没有人会因使人们遭受暴行而责备竞赛的监督。伤害诉讼在赛场之外。但是，就那些人涉及淤青、血污和肿胀的程度而言，他必会为他们设计冠冕、荣耀、礼物、政治特权、市民捐献、肖像、雕像，以及——世界所能给予的那类——永恒的名声、藉被铭记而复活。拳击手自己并不抱怨感受到疼痛，因为他希望如此；冠冕合拢伤口，棕榈枝掩盖血迹：他为胜利而非为伤害所激动。你见到快乐的人，会认为他受伤了吗？但即便是失败者自己，也不会因他的不幸而责备竞赛的监督。难道天主将祂自己的技艺种类和规则公之于众，置于这世界的空旷场地，供人、天使和众权力观看，竟是不合宜的吗？——为考验肉体与精神的坚定和忍耐？——给这人棕榈枝，给那人荣誉，给另一人公民特权，给另一人报酬？——也拒绝一些人，并在惩罚后羞耻地驱逐他们？你竟敢指示天主，在何时、以何种方式、在何地，为祂自己的队伍（竞争者）设立考验，似乎审判者也不适合宣布初步判决。那么，如果祂显明信德是为了受苦殉道，不是为竞赛，而是为它本身的益处，难道它就不该有某些希望的储备吗？为了这储备的增加，它可能克制自己的欲望，抑制自己的意愿，以便努力向上攀登，因为那些履行地上职务的人也渴望升迁？否则，如果我们父的家中有许多住所，若不是为了与功绩的多样性相符，又当如何？一颗星如何与另一颗星在光荣上有别，若不是因它们光线的差异？\BibleRef{格前 15:41}但更进一步，若因此，些许光芒的增加也适于崇高的信德，那么这获得应属于某种需付出巨大努力、剧烈痛苦、磨难、死亡之类。但请考虑报答：当肉体和生命被付出——在人内没有比它们更宝贵的了，一者来自天主的手，另一者来自祂的气息——正是这些东西被付出以获取那由这些利益构成的利益；正是这些东西被花费，而它们也是可获得的；相同的东西既是代价，也是商品。天主也预见了人类处境中其他弱点——敌人的诡计、受造物的欺骗性面相、世界的网罗；信德，甚至在圣洗后，也会陷于危险；大多数人，在获得救恩后，将再次失落，因玷污了婚宴礼服，因未能为灯盏预备油——他们将是那种须在山上和林间被寻找，并扛在肩上背回的人。因此，祂设立了安慰的第二项供给，以及最后的救援方式——殉道的战斗和血洗——此后便无危险。关于分享这些的人的幸福，\Person{达味}说：\BibleQuote{过犯得赦免、罪恶得遮盖的人，是有福的；主不归罪于他的人，是有福的。}因为，严格说来，不能再记算任何罪债反对殉道者，他们在（血）洗中放下了生命本身。因此，\BibleQuote{爱德遮盖许多罪过；}\BibleRef{伯前 4:8}并且爱天主，即尽全部力量（以这力量在忍受殉道时维持战斗），尽全部生命（为天主放下生命）\BibleRef{玛 22:37}，这使人为殉道者。你称这些为医治、劝告、审判方法、景象，甚至天主的野蛮吗？天主渴望人的血吗？然而我敢断然地说，祂渴望，如果人也渴望天国，如果人渴望确实的救恩，如果人也渴望第二次新生。这交换不会使任何人不满，因为它可以为自己辩护说，进行交换的双方分享了利益或伤害。\par

% structure: p0015 source=h2 target=section reason=relative-heading-rank; base=h2

\section{第七章}

如果那蝎子在空中甩动尾巴，仍以有一位杀人者为我们的天主来责备我们，我便要因从他异端之口发出的那股全然污秽的亵渎臭气而战栗；但我仍要拥抱这样一位天主，怀着由理性所生的确信，正是藉这理性，祂本人也借着祂自己\OriginalTerm{智慧}{Sophia}的位格，通过撒罗满的口，宣告自己比杀人者更甚：\OriginalTerm{智慧}{Sophia}，祂说，已杀死了她自己的儿女。\OriginalTerm{智慧}{Sophia}就是智慧。她确实智慧地杀死了他们，只要是为进入生命；也合理地杀死了他们，只要是为进入荣耀。论及为父者的杀人，哦，何等巧妙的形式！哦，罪行的机巧！哦，残忍的证明，祂杀人正是为此：被杀的人不致死亡！那么结果如何？智慧在出口之处受赞美，在歌咏中受颂扬；因为殉道者的死亡也受到赞颂。智慧在街市上坚定地行事，因为她杀死自己的儿子们带来了好的结果。不但如此，她在城墙顶上自信地发言，正如依撒意亚所说，这人大声说：\BibleQuote{我属于天主}，那人呼喊：\BibleQuote{因雅各伯的名}，另一人写着：\BibleQuote{因以色列的名}。\BibleRef{依 44:5}哦，良善的母亲！我自己也愿被列在她儿子之中，好让我被她所杀；我愿被杀死，好使我成为儿子。但她仅仅杀死她的儿子们吗？还是也折磨他们？因为我听见天主在另一处说：\BibleQuote{我必煅烧他们，如同烧炼金子；试炼他们，如同试炼银子。}\BibleRef{匝 13:9}诚然，藉着火与刑罚所提供的折磨，藉着信德的殉道考验。宗徒也知道他向我们描述了怎样一位天主，他写道：\BibleQuote{天主既然没有怜惜自己的儿子，反而为我们众人把他交出来，岂不也把一切和他一同赐给我们吗？}\BibleRef{罗 8:32}你看，天主的智慧甚至杀死了她自己的、首生的独生子，他必然要活着，不但活着，更要把其他的人也带入生命。我敢与天主的智慧同说：\BibleQuote{是基督为我们的过犯被交付。}\BibleRef{罗 4:25}智慧已经连自己也屠宰了。话语的特征不仅在于声音，也在于意义；它们不仅要用耳朵听，更要用心灵听。不明白的人，就相信天主是残忍的；虽然对于那不明白的人，也已经发出宣告，以约束他那非\Emph{正确}理解的严厉。宗徒说：\BibleQuote{因为谁曾知道主的心意？或者谁曾作过祂的参谋，教训祂？或者谁曾给祂指出理解之路？}\BibleRef{罗 11:34}然而，这个世界曾认为，以人祭来平息\Person{斯库提亚人}的\Person{狄安娜}、\Person{高卢人}的\Person{墨丘利}、\Person{阿非利加人}的\Person{萨图尔努斯}是合法的；在\Place{拉丁姆}，直到今日，\Person{朱庇特}还在城中心被人献上人血供他品尝；没有人对此加以议论，或认为这些祭献的发生没有理由，或认为这是按他的神的旨意发生的，而无价值。如果我们的天主，为了拥有祂自己的祭献，而要求为主殉道，谁又会因那致命的宗教、哀悼的仪式、祭台上的火堆、以及殡葬司祭而责备祂呢？谁不会反倒称那被天主吞噬的人为有福呢？\par

% structure: p0017 source=h2 target=section reason=relative-heading-rank; base=h2

\section{第八章}

因此，我们坚持这一个立场，并且仅就这个问题，召集一场辩论：\OriginalTerm{殉道}{martyrdom}是否由天主所命令？好让你相信它们是有理由被命令的，如果你知道是由祂所命令的，因为天主不会无故命令任何事。因为祂的圣人们的死亡在祂眼中是宝贵的，正如\Person{达味}所歌咏的，我想，那不是一般人常见的死亡，也不是众生所欠的债（那死亡甚至因罪恶而可耻，并且应追溯到\Emph{应受的惩罚}），而是\Emph{另一种}死亡，即在这项工作中所遇到的——为宗教作见证，并为了\OriginalTerm{义德}{righteousness}和\OriginalTerm{圣事}{sacrament}而坚持\OriginalTerm{宣认}{confession}的战斗。正如\Person{依撒意亚}所说：\BibleQuote{请看义人怎样灭亡，没有人放在心上；正直的人被除去，没有人理会：因为义人从邪恶面前灭亡，他将在埋葬时得到尊荣。}这里，你也既得到殉道的宣告，又得到它们所带来\Emph{的}报酬的宣告。实际上，从起初，义德就遭受暴力。一旦天主开始受到敬拜，宗教就得到敌意作为她的份。那曾取悦天主的人被杀害，而且是被他的兄弟杀害。从亲属的血开始，以便更轻易地去追寻陌生人的血，不敬神者把义人作为追逐的对象，最终，不仅是义人，甚至包括先知。\Person{达味}受迫害；\Person{厄里亚}被迫逃亡；\Person{耶肋米亚}被石头砸死；\Person{依撒意亚}被锯成两半；\Person{匝加利亚}在祭台与圣殿之间被屠杀，使坚硬的石头留下他血迹的持久痕迹。\BibleRef{玛 14:3}那人本人，在法律和先知书的末尾，被称为不是先知，而是使者，遭受耻辱的死亡，被斩首作为给跳舞女孩的奖赏。而且，那些惯于被天主圣神引导的人，曾被祂亲自引导去殉道；因此他们甚至已经要忍受他们曾宣告必须承受的。因此，那三个人的兄弟团体，当王像的奉献成为迫使市民献祭的场合时，深知\OriginalTerm{信德}{faith}（唯独在他们中没有被俘虏）要求什么——即他们必须抵制\OriginalTerm{偶像崇拜}{idolatry}至死。\BibleRef{达 3:12}因为他们也记起\Person{耶肋米亚}写给那些即将被掳的人的话：\BibleQuote{现在你们将看见巴比伦人的神祇被扛在肩上，是金、银、木制的，使外邦人恐惧。因此，你们要谨慎，不要完全像外邦人一样，当你们看见人群在前面和后面崇拜那些神祇时，不要被恐惧抓住，反而要在心里说：我们的本分是敬拜祢，主。}\BibleRef{巴 6:3}因此，他们从天主那里获得了信心，当他们以坚强的意志藐视国王对不顺从者的威胁时，说：\BibleQuote{我们没有必要回答你的这个命令。因为我们所敬拜的天主能够将我们从火窑中和你的手中拯救出来；然后就会向你们显明，我们既不侍奉你的偶像，也不崇拜你所设立的金像。}\BibleRef{达 3:16}哦，即使没有痛苦的殉道也是完美的！他们受了足够的苦！他们被烧得够了，天主因此保护了他们，以免他们似乎错误地代表了祂的大能。因为如果\Person{达理阿}关于天主的正当期望被证明是虚假的，那么狮子立刻当然会以其惯常的凶猛吞噬\Person{达尼尔}——他只敬拜天主，因此被\Person{加色丁人}控告和索要。至于其余的人，凡是天主的宣讲者，以及每一位敬拜者，比如那些被召唤去事奉偶像崇拜而拒绝服从的人，都应该受苦，也符合那个论点的要旨，即\OriginalTerm{真理}{truth}应被推荐给当时活着的人和后来的人——（即）真理辩护者自己的受苦本身为真理赢得了信任，因为除了拥有真理的人，没有人会愿意被杀。这样的命令以及实例，追溯到最早的时代，表明信徒有义务接受殉道。\par

% structure: p0019 source=h2 target=section reason=relative-heading-rank; base=h2

\section{第九章}

为防止古时可能曾有属于他们自己的圣事，我们必须审视现代基督宗教的体制，就好像它也来自天主，因而与\Emph{先前的}不同，并且在其行为准则上也与之对立，以至于其智慧竟不知杀害自己的孩子！显然，在基督身上，神性、旨意和教派都与\Emph{先前所知的}任何事物不同！祂要么从未命令过殉道，要么命令那些必须以不同于寻常意义来理解的殉道，且祂是这样一位：并不鼓动任何人冒这种风险，也不许诺给那些为祂受苦的人任何赏报，因为祂不愿他们受苦；因此，在颁布首要诫命时，祂说：\BibleQuote{为义而受迫害的人是有福的，因为天国是他们的。}以下这段话，\Emph{首先}适用于所有人，无有例外，然后特指宗徒们：\BibleQuote{人若因我的缘故辱骂你们、迫害你们，捏造各样坏话毁谤你们，你们是有福的。你们欢喜踊跃吧！因为你们在天上的赏报是丰厚的；因为他们的祖先也这样对待过先知。}因此祂同样预言了，他们自己也将如先知一样被杀害。即使祂曾规定，唯有当时那些宗徒在遵命的情况下要经历这一切迫害，那么，通过他们，连同完整的圣事、圣名的枝条、圣神的洗礼，关于忍受迫害的规条也必然指向我们，如同继承的门徒和宗徒种子的枝条。因为祂再次这样对宗徒们说训导的话：\BibleQuote{看，我派遣你们好像羊进入狼群中；}并说：\BibleQuote{你们要提防人，因为他们会把你们交给公议会，也会在会堂里鞭打你们；并且你们要为了我的缘故被带到总督和君王面前，对他们和外邦人作证，}等等。\BibleRef{玛 10:16}当祂加上\BibleQuote{兄弟要把兄弟送到死地，父亲要把孩子送到死地；儿女要起来反抗父母，害死他们}时，祂清楚地宣布，对于其他人，也将遭受这种不义行为，而我们在宗徒身上并未见到\Emph{实例}。因为他们当中没有一个人经历过父亲或兄弟的出卖，而我们的许多人却有过这样的经历。然后祂又转向宗徒：\BibleQuote{你们要为我的名字被众人憎恨。}我们更是如此，因为我们还要被父母出卖！因此，祂将这种背叛——时而分配给宗徒，时而分配给众人——将同样的毁灭倾注在所有拥有这名号的人身上，那名号连同其成为憎恨对象的条件，将常在他们身上。\BibleQuote{但那坚持到底的，终必得救。}坚持什么？无非是迫害、背叛、死亡。因为坚持到底无非就是受苦到底。因此紧接着是：\BibleQuote{门徒不能超过师傅，仆人也不能超过主人。}因为，既然师傅和主自己也在忍受迫害、背叛和死亡中坚定不移，那么祂的仆人和门徒就更应该忍受同样的遭遇，免得他们看起来胜过祂，或免于不义的攻击；因为对他们而言，足以荣耀的正是与主和师傅的苦难相似。祂为了预备他们忍受这些，提醒他们不要惧怕那些只能杀死身体，却不能毁灭灵魂的人；而应该敬畏那位有能力同时杀死身体和灵魂，并将他们投入地狱中的。请问，这些只能杀死身体的是谁？不是前面提到的总督和君王吗？我想，是凡人吧？那又能掌管灵魂的是谁？不是唯独天主吗？祂岂不就是那位威胁\Emph{将来}烈火的主？若没有祂的旨意，连一只麻雀也不会掉在地上——也就是说，人的两种实体——肉身或灵魂——一个也不会，因为我们的头发也都被祂数过了。因此，不要惧怕。当祂加上\BibleQuote{你们比许多麻雀还贵重}时，祂应许我们不会徒然——即不是毫无益处——掉在地上，如果我们选择被人杀死而不是被天主杀死。\BibleQuote{所以，凡在人面前承认我的，我在我天上的圣父面前也必承认他；凡在人面前否认我的，我在我天上的圣父面前也必否认他。}我认为，在宣布承认和否认时所用的措辞以及解释方式都是清楚的，尽管表达方式不同。承认自己是基督徒的人，就见证他是属于基督的；属于基督的人，就是在基督内。若他在基督内，他承认时必定是在基督内承认，因为他承认自己是基督徒。因为他不可能不在基督内而成为基督徒。此外，在基督内承认时，他也承认了基督：因为作为基督徒，他在基督内，同时\Emph{基督}本身也在他内。若是你提到白天，你同时也展示给我们白天的光体——太阳，尽管你可能没有提到光。因此，虽然祂没有明确说\BibleQuote{谁承认我}，但日常生活承认的行为与主的宣言并非不同。因为承认自己是什么，即基督徒，也就承认了那使他成为此身份者，即基督。因此，否认自己是基督徒的人，就在基督内否认了，因为他否认自己在基督内，同时否认了自己是基督徒；另一方面，否认基督在他内，就是否认他在基督内，从而也否认了基督。所以，凡在基督内否认的，就否认了基督；凡在基督内承认的，就承认了基督。因此，即使主只是宣告了关于承认的教导，也就足够了。因为从祂提出承认的方式，也能预先判定其反面——即否认——的景况：否认会被主以否认来回应，正如承认以承认来回应。因此，既然在承认的形态中也可以看出否认的状态，那么显然，主以不同于谈论承认的措辞所宣告的\BibleQuote{谁否认我}（而不是\BibleQuote{谁在我内否认}），是属于另一种方式的否认。因为祂预见，这种暴力形式也往往紧随其后，即当有人被迫放弃基督徒名号时，那否认自己是基督徒的人，将被强迫以亵渎祂的方式否认基督本身。唉，不久之前，我们曾战栗地目睹一些人在此方式下与他们那曾有吉兆的完整信仰挣扎。因此，若有人说：\Quote{我虽否认自己是基督徒，却不会被基督否认，因为我并未否认祂本身。}这是徒劳的。因为从那个否认中，既然他因否认自己是基督徒而否认了基督在他内，也就\Emph{同样}否认了基督本身。不仅如此，因为祂也用羞耻来威胁羞耻：\BibleQuote{凡在人面前以我为耻的，我在我天上的圣父面前也必以他为耻。}因为祂知道，否认往往最由羞耻产生，内心的状态显露在额头上，羞耻的创伤先于身体的创伤。\par

% structure: p0021 source=h2 target=section reason=relative-heading-rank; base=h2

\section{第十章}

至于那些认为宣认并非在此——即在此大地环境内、在此生存时期中、在具有我们人人所共有的本性的世人面前——被规定的人，他们的假设何其荒谬，与我们在这些土地上、在此生命中、在人间权威之下所经验的一切秩序相悖！\newline{}无疑地，当灵魂离开躯体，开始在天上的各层被审判，关乎他们来到耶稣面前时所担负的盟约，并就异端者的那些隐藏奥秘受审问时，他们就必须在那真实的权势和真实的人面前宣认——即瓦伦提努斯（Valentinus）的特莱提（Teleti）、阿巴斯坎提（Abascanti）和阿西内提（Acineti）！\newline{}因为，\Emph{他们}说，就连德谬哥（Demiurge）自己也没有一致地赞同我们这世界的人，他看他们如桶中的一滴水\BibleRef{依 40:15}，如禾场上的灰尘，如唾沫和蝗虫，\Emph{甚至}与禽兽同列。显然，这样写着。然而，我们不应因此就以为除我们之外还有另一种人，这种人——在所述情况下显然\Emph{如此}——能够在不使\Emph{两种类型之间的}比较失效的情况下，兼具种族的特征和独特的属性。因为，即使生命被玷污，以致被定罪而受轻蔑，可与被轻蔑之物相比，但本性并未立即被剥夺，以致在其名下可假定有另一种人。相反，虽然生命蒙羞，本性却被保全；基督所认识的人，正是那些祂说\BibleQuote{人们说我是谁？}\BibleRef{玛 16:13}以及\Quote{你们愿意人怎样待你们，你们也要怎样待人。}所涉及的人。请思考，祂是否保全了这样一类人，以致祂从他们身上寻求对自己的见证，并且由祂命令他们彼此施行公平对待的人所组成？\newline{}但如果我急切地要求那些天上的人被描述给我听，阿拉托斯（Aratus）更容易在星座间勾勒出珀耳修斯（Perseus）、刻甫斯（Cepheus）、厄里戈涅（Erigone）和阿里阿德涅（Ariadne）。然而，是谁阻止了主清楚规定，宣认也必须在人前进行，正如祂明确宣告祂自己的人要在哪里宣认一样？以至于这句话本可以这样表述：凡在人前在天上承认我的，我在我天上的父面前也必承认他。祂本应将我从这关于地上宣认的错误中拯救出来，因为祂若命令了在天上的宣认，就不愿我参与地上的宣认；因为除了地上的居民，我不认识别的人，甚至到那时为止也没有人在天上被看见。\newline{}此外，所断言的事情有何可信之处？死后被提升到天上，在那里受考验——而我未受考验就不会被转移到那里——在那里就一命令受审，而我要不先进入那地方就无法到达？天堂对于基督徒是在其道路之前就敞开的；因为除了天堂向他敞开的人，没有通往天堂的道路；而他到达了就进入。\newline{}我听见你断言在天门口守卫的权势，依据罗马人的迷信，有某个卡努斯（Carnus）、福尔库鲁斯（Forculus）和利门提努斯（Limentinus）？你在栏杆处安排了哪些权势？\newline{}如果你曾在达味（David）的《圣咏》中读到：\BibleQuote{众城门哪，你们要抬起头来！永久的门户，你们要被抬起！荣耀的君王要进来！}如果你也从亚毛斯（Amos）听见：\BibleQuote{那在天上建造楼阁，在地上奠定根基，召唤海水，使其倾泻在地面上的——那一位}\BibleRef{亚 9:6}。要知道，那上升之路此后被主的脚步夷为平地，入口此后靠基督的大能打开，基督徒在门槛上不会遇到延误或审查，因为他们在那里不是要被辨别彼此，而是被认领，不是被审问，而是被接纳。\newline{}因为，即使你认为天堂仍然关闭，但请记住，主将钥匙留给了伯多禄（Peter），并通过他留给了教会；每一个在此地受审并宣认的人，都将带着这钥匙。但魔鬼顽固地声称我们必须在天上宣认，为要说服我们在地上否认。\newline{}我将预先送去精美的文件，确实，我将携带精美的钥匙——就是那只能杀身体却不能伤害灵魂者的恐惧。我将因忽视这命令而蒙恩；我将在天上光荣站立，而我在地上站立不住；我将向更大的权势持守，而我向较小的权势屈服了；我将最终配得被接纳，而此刻被关在门外。\newline{}很容易想到进一步评论：\Quote{如果必须在天上宣认，那么也必须在本地否认。}因为二者在同一个地方；对立面总是相伴相随。甚至需要在天上进行迫害，这是宣认或否认的场合。\newline{}那么，你这最傲慢的异端者，为何不把威吓基督徒的整套手段，尤其是那名字本身的仇恨，转移到上面的世界？那里基督坐在圣父的右边。你将在那里设立犹太人的会堂——迫害的泉源——宗徒们在其前受鞭打，还有外邦人的集会连同它们的马戏团，他们也随时呼喊：\Quote{处死那第三种族！}\newline{}但你必须在同一地方也带出我们的弟兄、父亲、儿女、岳母、儿媳和我们的家人，通过他们出卖已被指定；同样还有君王、总督和武装权威，在他们面前争讼的问题必须处理。肯定地，天上也将有一座监狱，缺乏阳光或充满忘恩负义的光，也许还有锁链，以及作为刑架的、旋转\Emph{苍穹}的轴本身。\newline{}然后，若基督徒要被石头砸死，冰雹会近在眼前；若被焚烧，雷电就在手边；若被斩杀，武装的俄里翁（Orion）将履行职责；若被野兽致死，北方会派出熊，黄道带会派出公牛和狮子。那忍耐这些攻击到底的，必要得救。\newline{}那么，在天上将有终结、苦难、杀戮和初次宣认吗？而且，这一切所需的肉体在哪里？那只能被人杀死的身体在哪里？\newline{}无误的理性命令我们以近乎戏谑的方式陈述这些事；没有人能推开这异议的障碍，而不被迫将迫害的整套手段、所有为此目的而提供的强大工具，转移到他所放置宣认法庭的地方。既然宣认由迫害引出，而迫害以宣认告终，那么决定其入口和出口（即开始与结束）的工具必然同时伴随它们。\newline{}但是，对名字的仇恨将在这里，迫害在这里爆发，出卖在这里引人，审问在这里施压，酷刑在这里肆虐，而宣认或否认在此地上完成这全套程序。因此，若其他事物在此，宣认也不在别处；若宣认在别处，其他事物也不在此。当然，其他事物不在别处，所以宣认也不在天上。\newline{}或者，若他们坚持天上的审问和宣认方式是另一种，那也必然要发明一种非常不同的、与圣经所指示的方式相对立的程序。我们可以说：让他们思考（他们所想象的是否存在），如果这属于地上审问和宣认的程序——它以迫害为起源，以国家中的纷争为理由——被保持在自己的信德中，如果我们必须正如经上所记载的那样相信，正如所说的话那样理解。\newline{}我在此忍受全部过程，主自己并没有指定另一个世界\Emph{给我这样做}。因为，在论及宣认和否认之后，祂又加了什么？\BibleQuote{你们不要以为我来是给地上送和平，而是送刀剑}——无疑是在地上。\BibleQuote{因为我来是要使儿子反对他的父亲，女儿反对她的母亲，儿媳反对她的婆婆。人的仇敌就是自己家里的人。}\BibleRef{玛 10:34}因为这样就实现了：兄弟把兄弟交付死地，父亲把儿子交付死地；儿女起来反对父母，并害死他们。\BibleQuote{那忍耐到底的，那人必要得救}\BibleRef{玛 10:21}。因此，这属于主刀剑的整个程序——它不是被送向天上，而是被送向地上——也就使宣认也在那里，宣认要以忍耐到底并死于苦难而完成。\par

% structure: p0023 source=h2 target=section reason=relative-heading-rank; base=h2

\section{第十一章}

同样，我们坚持认为，其他的宣告也指向殉道的处境。耶稣说：\Quote{谁若重视自己的生命甚于我，就不配是我的门徒}\BibleRef{路 14:26}——也就是说，谁若宁愿背弃我而活着，也不愿承认我而死去；又说：\Quote{谁获得自己的生命，必要丧失生命；谁为我的缘故丧失自己的生命，必要获得生命}\BibleRef{玛 10:39}。因此，的确，那在赢得生命时背弃信仰的，就丧失生命；而那以为通过背弃赢得生命的人，却要在地狱中丧失生命。相反，那因承认信仰而被杀的人，暂时丧失生命，但必将获得永生。总之，当权者自己，当他们劝人背弃信仰时，常说：\Quote{救你的命吧！}和\Quote{不要丧失你的生命！}基督说话时，怎能不吻合基督徒将受到的对待呢？但当祂禁止在审判席前思虑如何回答时\BibleRef{玛 10:19}，祂是在预备祂的仆人去\Emph{面对等待他们的事}，祂保证圣神将\Emph{藉着他们}说话；当祂希望弟兄在狱中被探访时\BibleRef{玛 25:36}，祂是在命令人们关怀那些即将承认信仰的人；当祂断言天主必为祂的选民伸冤时，祂是在安慰他们的痛苦\BibleRef{路 18:7}。在种子萌芽后枯萎的比喻中\BibleRef{玛 13:3}，祂描绘了一幅关于逼迫热火的图画。如果这些宣告不被按照其原意理解，那无疑它们所指的与字面所示的不同；那么言词是一回事，含义是另一回事，就如寓言、比喻、谜语一样。因此，无论这些蝎子\EditorialNote{指异端者}可能捕捉到什么推理之风，无论他们以何种诡计攻击，现在有一条防线：我们将诉诸事实本身，看它们是否如圣经所预言的那样发生；因为如果在实际事实中找不到那（似乎）相同的事，那么圣经中便另有所指。因为所写的必须应验。此外，所写的将要应验，除非情况不同。但是，看！我们因这名而成为众人憎恨的对象，正如所写的；我们也被至亲出卖，正如所写的；被带到官长面前，受审问，受折磨，承认信仰，被残忍杀害，正如所写的。主就是这样预定的。如果祂预定这些事以别的方式发生，为什么它们不以祂所预定的方式发生，即照祂所预定的那样？但它们的发生并不超出祂的预定。因此，它们如何发生，祂就如何预定；祂如何预定，它们就如何发生。因为既不会允许它们超出祂的预定而发生，祂也不会预定与祂意愿相悖的事。所以，这些经文的意义不会与我们在实际事实中所认识的不同；或者，如果所宣告的事件尚未发生，那么那些未宣告的事件如何正在发生？因为这些正在发生的事件并未被宣告，如果所宣告的是别的，而不是这些正在发生的事件。那么，既然在现实生活中遇到的事件被认为在言词中有不同的含义，那如果发现它们以\Emph{不同于所启示}的方式发生，又会怎样？但这将是信仰的乖僻：不相信已被证明的事，却假设未被证明的事是真的。对这种乖僻，我还要提出以下反对：如果这些正如所写发生的事件不是所宣告的那些，那么（所指的）那些也不应按所写的方式发生，以免它们也因这些\Emph{其他}事件之例而有被排除的危险，因为言词与事实不符；结果，所宣告的事件在发生时也未被看见，如果它们的宣告与它们必须发生的方式不符。那么，那些未按发生方式被宣告的事件，如何能被相信（已经发生）？因此，异端者因不相信按照已被显示发生的方式所宣告的，反而相信连宣告都未曾有的事。\par

% structure: p0025 source=h2 target=section reason=relative-heading-rank; base=h2

\section{第十二章}

如今，谁能比基督自己的学派更明白圣经的精髓呢？——就是那些主亲自拣选为门徒，必然要在一切事上受全备教导，并指派给我们作师傅、在一切事上教导我们的人。祂更愿意向谁揭示祂自己言语的隐藏含义，而不是向那些祂曾显示自己荣耀相似的人——即\Person{伯多禄}、\Person{若望}和\Person{雅各伯}，以及后来的\Person{保禄}，祂在保禄殉道之前还赐予他参与（乐园的）喜乐？或者他们也写下与自己想法不同的话——师傅用欺骗，而不用真理？\Person{伯多禄}在致本都基督徒的信中，无论如何说：\BibleQuote{如果你们耐心受苦，不因作恶而受罚，这确实是多大的光荣！因为这是可爱的特质，\Emph{而且}你们正是为此蒙召，因为基督也为我们受苦，给你们留下了榜样，叫你们追随祂的脚步。} \BibleRef{伯前 2:20} 又说：\BibleQuote{亲爱的，你们中间的火炼的试验不要以为奇怪，似乎是遭遇非常的事；倒要欢喜，因为你们是与基督一同受苦，使你们在祂荣耀显现的时候，也可以欢喜快乐。你们若为基督的名受辱骂，便是有福的，因为荣耀和天主的圣神安息在你们身上；只是你们中间不可有人因为杀人、偷窃、作恶、好管闲事而受苦。但若作为基督徒受苦，却不要羞耻，倒要因这名荣耀天主。} \BibleRef{伯前 4:12} 事实上，\Person{若望}劝勉我们甚至要为弟兄舍命，\BibleRef{若壹 3:16} 断言爱里没有恐惧：\BibleQuote{因为完美的爱驱除恐惧，因为恐惧含有惩罚；那恐惧的人在爱里并不完美。} \BibleRef{若壹 4:18} 我们最好将（此处所指的）恐惧理解为哪一种？岂非那导致否认的恐惧？他断言哪一种爱是完美的？岂非那驱走恐惧、赐予承认勇气的爱？他指定什么作为恐惧的惩罚？岂非那否认者将要付上的，即灵魂和身体都要在地狱里被消灭？并且如果他教导我们必须为弟兄死，何况为主呢？——他自己也藉着自己的默示录充分准备好给出这样的劝告！因为圣神曾向\Place{斯米纳}教会的天使传达命令：\BibleQuote{看，魔鬼要把你们中间几个人下在监里，叫你们被试炼十日。你务要至死忠心，我就赐给你那生命的冠冕。} \BibleRef{默 2:10} 也给\Place{培尔加摩}教会的天使（提及了）\Person{安提帕}，\BibleRef{默 2:13} 那位至死忠心的殉道者，在撒殚居住之处被杀。也给\Place{非拉德非雅}教会的天使（表明）\BibleRef{默 3:10}，那没有否认主之名的人得以脱离最后的试炼。接着，圣神向每位得胜者应许：现在有生命树，并免于第二次死亡；现在有隐藏的玛纳和白亮的石头，以及那（除了领受者以外无人知道的）新名；现在有以铁杖治理的权柄，和晨星的明亮；现在有穿白衣，名字不从生命册上涂抹，并在天主的殿中作柱子，上面刻着天主和主的名，以及天上的耶路撒冷；现在有与主同坐在祂的宝座上——这原是\Person{载伯德}的儿子们曾坚决要求的。\BibleRef{玛 20:20} 请问，这些如此蒙福的得胜者是谁？不就是严格意义上的殉道者吗？因为胜利属于那些战斗者；而战斗属于那些流血者。然而殉道者的灵魂暂时安息在祭台下，\BibleRef{默 6:9} 并藉着确切的复仇希望支撑他们的忍耐；他们穿着白衣，戴着灿烂的光环，直到其他人也完全分享他们的荣耀。因为又有一大群无数的人显现出来，穿着白衣，手持胜利的棕榈枝，无疑是在庆祝对假基督的胜利，因为一位长老说：\BibleQuote{这些人是从大患难中出来的，曾用羔羊的血把衣裳洗白净了。} \BibleRef{默 7:14} 因为肉身是灵魂的衣裳。不洁固然藉着圣洗被洗去，但污点却藉着殉道变为灿烂的白。因为\Person{依撒意亚}也预言，从红色和朱红中要生出雪和羊毛的白。\BibleRef{依 1:18} 同样，当大巴比伦被描绘成喝醉了圣徒的血，\BibleRef{默 17:6} 无疑她醉酒所需的供应是由殉道的杯提供的；同样，对殉道的恐惧会带来什么痛苦，也照样显示出来。因为在所有被弃绝的人中，不，其中居首的，是那些胆怯的。\BibleQuote{但胆怯的，} \Person{若望}说——接着是其他——\BibleQuote{将要在火与硫磺的湖里有分。} \BibleRef{默 21:8} 因此，恐惧——如他书信所说，爱驱除恐惧——含有惩罚。\par

% structure: p0027 source=h2 target=section reason=relative-heading-rank; base=h2

\section{第十三章}

然而，宗徒\Person{保禄}，曾为迫害者，首先流了教会之血，但后来他将刀剑换成笔，将匕首转为犁，原是本雅明的一只豺狼，后来却如\Person{雅各伯}般亲自供应粮食——他（我说）如何为殉道辩护，他自己也当选择殉道，当他为得撒洛尼人而欢喜时，说：\BibleQuote{所以我们在天主的各教会中为你们夸耀，因为你们在所受的一切迫害与磨难中，显示了坚忍和信德；这正是天主公义审判的明证，使你们堪当祂的国，你们也正是为这国而受苦！}\BibleRef{得后 1:4} 同样在致罗马人书中也说：\BibleQuote{不但如此，我们也在磨难中夸耀，因为知道磨难生坚忍，坚忍生老练，老练生望德，望德不使人蒙羞。}\BibleRef{罗 5:3} 又说：\BibleQuote{若是子女，便是承继者，是天主的承继者，与基督同为承继者：只要我们与祂一同受苦，也必与祂一同受光荣。因为我算定，现时的苦难，与将来在我们身上要显示的光荣，是无法相比的。}\BibleRef{罗 8:17} 因此他随后说：\BibleQuote{谁能使我们与天主的爱隔绝呢？是磨难吗？是困苦吗？是迫害吗？是饥荒吗？是赤身露体吗？是危险吗？是刀剑吗？正如经上记载：\InnerQuote{为了祢的缘故，我们终日被置于死地，被视作待宰的羊。}然而，靠着那爱我们的主，我们在这一切事上，大获全胜。因为我深信：无论是死，是生，是权势，是高天，是深渊，或是任何别的受造物，都不能使我们与天主的爱隔绝，这爱是在我们的主基督耶稣内的。}\BibleRef{罗 8:35} 再者，他在向格林多人叙述自己的苦难时，确已断定必须忍受苦难：\BibleQuote{（他说）论劳苦，更为频繁；论监禁，更是无数；论死亡，屡遭危险。被犹太人鞭打五次，每次四十下减去一下；被棍击三次；被石头砸一次；}\BibleRef{格后 11:23} 等等。若这些酷刑似乎比殉道更沉重，但他又再说到：\BibleQuote{因此，我为基督的缘故，在软弱、凌辱、急难、迫害、困苦中，感到喜乐。}\BibleRef{格后 12:10} 他在书信前段经句中还说：\BibleQuote{我们的处境是：处处受困，却不被窘迫；有所缺乏，却不至匮乏；受迫害，却不被遗弃；被打倒，却不至灭亡；身上常带着耶稣的死。}\BibleRef{格后 4:8} \BibleQuote{但}他说，\BibleQuote{我们外在的人虽在毁坏}——无疑，指肉身因迫害之力——\BibleQuote{内心的人却日日更新}——无疑，指灵魂藉对预许之望德。\BibleQuote{因为我们这暂时的轻微苦难，为我们成就卓越而永恒的荣耀分量；我们并不注目看得见的事物，而是注目看不见的事物；因为看得见的是暂时的}——他指的是苦难——\BibleQuote{看不见的却是永恒的}——他是在预许赏报。他在狱中写信给得撒洛尼人时，确曾肯定他们是有福的，因为不仅赐予他们信基督，也为祂受苦。\BibleQuote{（他说）你们有同样的斗争，正如你们在我身上见过，如今又在我身上听见的。}\BibleRef{斐 2:17} \BibleQuote{即使我作为奠祭被浇奠在祭献上，我也与你们众人一同喜乐，一同欢庆；你们也要同样与我一同喜乐，一同欢庆。}你看他如何断定殉道的福乐，为尊崇这福乐，他预备了彼此喜乐的庆节。当他最终临近所渴慕之境，满心欢喜地展望前面时，他这样写给\Person{弟茂德}：\BibleQuote{因为我已被奠祭，我离世的时刻到了。那美好的仗我已打过，赛程我已跑完，信德我已持守；从此以后，正义的冠冕已为我预备，主在那一天要赐给我——}\BibleRef{弟后 4:6} 无疑，指苦难的冠冕。他在先前段落中也曾给予足够劝勉：\BibleQuote{这话是确实的：如果我们与基督同死，也必与祂同生；如果我们受苦，也必与祂一同为王；如果我们否认祂，祂也必否认我们；如果我们不信，祂仍是信实的——祂不能否认自己。}\BibleRef{弟后 2:11} \BibleQuote{所以，不要以我们主的见证为羞耻，也不要以我这为祂被囚的人为羞耻；}\BibleRef{弟后 1:8} 因为他先前说过：\BibleQuote{因为天主没有赐给我们恐惧之灵，而是刚强、仁爱、谨慎之灵。}\BibleRef{弟后 1:7} 因为我们是以对天主的爱之刚强，并以谨慎之心受苦，当我们因自己的清白而受苦时。再者，若祂在何处命令人忍受，那么祂为苦难预备这不正是为何？若祂在何处引人脱离偶像崇拜，那么殉道岂非首当其冲，引人离去而受害？\par

% structure: p0029 source=h2 target=section reason=relative-heading-rank; base=h2

\section{第十四章}

无疑，宗徒劝告罗马人\BibleRef{罗 13:1}要服从一切权柄，因为没有权柄不是出于天主，并且（掌权者）不是无故带剑，他是天主的仆人，而且，他说，也是报复者，向作恶的人施行忿怒。因为他先前也曾这样说：\BibleQuote{因为统治者不是使善行恐怖，而是使恶行恐怖。你愿意不怕权柄吗？你当行善，就必得它的称赞；因为他是天主的仆人，是为了你的益处。你若作恶，就该惧怕。}因此，他命令你们服从权柄，不是为他自己逃避殉道提供机会，而是当他为善行辩护时，也视权柄为仿佛赐予义德的助手，神圣法庭的婢女，后者甚至在此世就预先对有罪者作出判决。接着他又指出他希望你们如何服从权柄，命令你们\BibleQuote{当纳贡的，给他纳贡；当完税的，给他完税}\BibleRef{罗 13:6}，也就是说，凯撒的归凯撒，天主的归天主\BibleRef{玛 22:21}；但人唯独是天主的财产。无疑，\Person{伯多禄}\BibleRef{伯前 2:13}也曾同样说，君王固然当受尊敬，但君王受尊敬，\Emph{惟独}当他守在自己的范围，远避僭取神圣尊荣时；因为父亲和母亲将与天主一同被爱，而不与祂等同。此外，甚至生命也不可被允许爱得超过天主。\par

% structure: p0031 source=h2 target=section reason=relative-heading-rank; base=h2

\section{第十五章}

那么，如今宗徒们的书信也众所周知。你们会说，难道\Emph{我们}这些在各样事上只是天真无邪的灵魂和鸽子的人，就喜欢误入歧途吗？我想是由于贪生怕死。但就算这样，意义离开了他们的书信。然而，我们知道宗徒们忍受了这样的苦难：教训是清楚的。当我通读\WorkTitle{宗徒大事录}时，我只看到这一点；我完全不寻求别的。那里的监狱、锁链、鞭打、大石、刀剑、犹太人的攻击、外邦人的集会、护民官的控告、君王的审案、总督的审判席和凯撒的名号，都不需要解释。\Person{伯多禄}被击打，\Person{斯德望}\Emph{被石头}打死\BibleRef{宗 7:59}，\Person{雅各伯}被杀如同祭台上的祭品，\Person{保禄}被斩首，都已用他们的血写就。若异端者希望他的信靠基于公共记录，帝国的档案将会作证，耶路撒冷的石头也会作证。我们阅读\WorkTitle{凯撒列传}：在罗马，\Person{尼禄}是第一个用血玷污新生信仰的人。然后\Person{伯多禄}被别人束上\BibleRef{若 21:18}，被钉在十字架上。然后\Person{保禄}获得与罗马公民身份相称的诞生，在罗马他因殉道而重生更显尊贵。无论我在哪里读到这些事件，一读我就学会受苦；对我来说，我以谁为殉道的老师并不重要，无论是宗徒的言论还是他们的死亡，只不过在他们的死亡中我也回忆他们的言论。因为他们不会遭受任何他们先前不知道必须遭受的事。当\Person{阿加波}也利用相应的行动预言锁链等候\Person{保禄}时，门徒们哭泣着恳求他不要冒险前往耶路撒冷，却是徒然。\BibleRef{宗 21:11}至于他，有意要阐明他一向教导的，他说：\BibleQuote{你们为什么哭泣，使我心碎？至于我，我不但愿意受捆绑，甚至愿意在耶路撒冷为我的主\Person{耶稣基督}的名而死。}于是他们退让说：\BibleQuote{愿主的旨意成就。}无疑，他们确知苦难包含在天主的旨意中。因为他们曾试图阻止他，意图不是劝阻，而是表示对他的爱；是渴求（保全）宗徒，而不是劝告反对殉道。即使那时有一个\Person{普罗迪科}或\Person{瓦伦提诺}站在旁边，建议人不在地上在人前宣认，而且事实上更不应该宣认，以免天主（似乎）渴求血，基督渴求救苦的回报，好像祂祈求这个，为使自己也能通过它获得救恩，他便会立刻从天主的仆人那里听到魔鬼从主那里听到的话：\BibleQuote{\Person{撒殚}，退到我后面去！你是我的绊脚石。经上记载：你要朝拜主，你的天主，惟独事奉祂。}但即使在现在，他听到这（斥责）也是应当的，因为他很久以后倾泻了这些毒药，尽管如此，这些毒药也不容易伤害任何一个软弱的人，只要有人有信心，在受伤之前，甚至之后立即喝下我们的这剂药。\par

\SourceRecord{Translated by S. Thelwall.；Ante-Nicene Fathers；Vol. 3.；Edited by Alexander Roberts, James Donaldson, and A. Cleveland Coxe.；Buffalo, NY: Christian Literature Publishing Co.,；1885.；<http://www.newadvent.org/fathers/0318.htm>.}
